% Introduction der Thesis, Entwurf 2026-09-12. Zwei Teile: (1) die Motivation
% des Delta-Kapitels, aus chapter_delta.tex \section{Motivation} hierher
% gezogen (Absaetze 1-4 dort; Gleichungen und Architektur bleiben im
% Kapitel); (2) das Setting des OMol25-Kapitels in der zuletzt vereinbarten
% Fassung. Labels: ch:background, ch:delta, ch:mr, sec:mace-training,
% sec:delta-architecture, sec:delta-data, sec:mr-methods, app:marks.

\chapter{Introduction}
\label{ch:introduction}

% ---------------------------------------------------------------- Teil 1
% Wortlaut aus chapter_delta.tex \section{Motivation}, Absaetze 1-4.
% Geaendert nur: Platzhalter aufgeloest, Praemisse benannt, Gleichungen und
% Architektur im Kapitel belassen, Roadmap angehaengt.

Machine-learned interatomic potentials (MLIPs) such as PaiNN or MACE
are commonly trained on DFT calculations. Trained on reaction paths,
they can search transition states and predict barriers at a fraction
of the cost of DFT. Transition1x~\cite{schreiner2022} was the first
large dataset built for this purpose: 9.6 million DFT energies and
forces on and around the reaction paths of 10\,073 organic reactions,
where earlier datasets held equilibrium structures and their
surroundings. It demonstrated that a PaiNN model trained on these
${\sim}10$ million reaction-path geometries reproduces DFT energies
and forces. However,
Transition1x was computed at \mbox{$\omega$B97X/6-31G(d)}, a
comparatively cheap level of theory chosen to make the generation of
ten million data points computationally feasible. More accurate
levels of theory exist. The range-separated
meta-GGA $\omega$B97M-V with a triple-$\zeta$ basis is among the most
accurate DFT methods for barrier heights~\cite{mardirossian2016,
goerigk2017}. Any model trained on the Transition1x dataset can only
be as accurate as the level of theory the dataset was generated with,
so upgrading the level of theory is desirable.

Such an upgrade does not necessarily require repeating the NEB searches at the
higher level of theory. A dataset like Transition1x consists of two
things: the geometries, which sample the configuration space along
reaction paths, and the energy and force labels attached to them. The
premise is that the geometries remain valid at any level of theory:
an MLIP only needs them to cover the relevant region, which the
Transition1x paths already do. Only the labels are tied to the cheap
level. Replacing them requires one single-point calculation per geometry at
the higher level; this is what relabelling means. It is far cheaper
than relaxing the paths again at that level, which would take hundreds
of gradient calculations per reaction instead of one per geometry
kept.

OMol25~\cite{omol25} relabelled the entire Transition1x dataset at
\mbox{$\omega$B97M-V/def2-TZVPD}\footnote{def2-TZVPD adds diffuse
functions to def2-TZVP; this work omits them for cost.} using
unrestricted Kohn--Sham\footnote{Chapter~\ref{ch:delta} uses restricted
Kohn--Sham throughout. The implications of unrestricted versus
restricted Kohn--Sham are discussed in Chapter~\ref{ch:mr}.}.

In Chapter~\ref{ch:delta}, we pursue a different goal. Relabelling a
whole dataset costs one calculation per geometry at the higher level,
and any way to reach that level at a fraction of the cost is
desirable. We therefore ask whether it can be reached by relabelling
only a small part of the data: about 80\,000 geometries, under 1\,\%
of the dataset. A MACE model is trained on Transition1x to predict the cheap level,
\mbox{$\omega$B97X/6-31G(d)} (Section~\ref{sec:mace-training}), and
a small correction head on top of it
(Section~\ref{sec:delta-architecture}) learns the difference between
the cheap and the expensive level, \mbox{$\omega$B97M-V/def2-TZVP}. This is feasible with so little data because
the difference is nearly constant within a reaction
(Section~\ref{sec:delta-data}): a smooth offset is far easier to learn
than the potential energy surface itself. Base model and correction
head combine into one corrected potential, MACE+$\Delta$, which can be
used wherever MACE can: for single-point predictions or as the
potential driving a NEB optimisation. The chapter evaluates it on 30
test reactions grouped by multireference (MR) character, at
fixed geometries and in searches the model drives itself. It finds that the force errors drop to a third on all reactions, and
that the correction reproduces the higher level where the
multireference character is low: there the barrier errors drop to a
third as well, while on the mid- and high-MR reactions the correction
overshoots.

% ---------------------------------------------------------------- Teil 2

Chapter~\ref{ch:delta} raises the level of theory of the labels and
leaves the geometries unchanged, on the premise that changing the labels is
enough to lift the model to the higher level of theory. That is how the field raises fidelity: delta
learning~\cite{ramakrishnan2015}, multi-fidelity
training~\cite{messerly2025}, dataset merging~\cite{shiota2024} and
reactive potentials with semi-empirical geometries~\cite{ren2026} all
supply labels at a higher level of theory on geometries they do not
re-optimise. OMol25~\cite{omol25} went one step further: it took
Transition1x, a dataset generated restricted at a cheap level of
theory, and relabelled the full set, not a fraction of it,
unrestricted at an expensive one, again without re-optimising a single
path. Two questions lie underneath this
practice, and this thesis asks both:

\noindent
\begin{minipage}{\linewidth}
\begin{description}
  \item[Research question 1 (RQ1)] Can a model be lifted to a higher
    level of theory by relabelling only a subset of its training data?
  \item[Research question 2 (RQ2)] Can a model learn a surface from
    paths that were never relaxed on that surface?
\end{description}
\end{minipage}

Chapter~\ref{ch:delta} answers RQ1. For that chapter, RQ2 needs no
test: both levels of theory are restricted, and when the transition
states are searched with DFT at either level, the two results lie
within 0.007\,\AA{} of each other. A path relaxed at the cheap level
therefore already covers the immediate region the transition state
lies in when calculated with the expensive method, and relabelling the
path without re-optimising it is enough.

This stops being true when the spin formalism changes as well. At a
broken-symmetry transition state the restricted solution is unstable
(Section~\ref{sec:background-stability}), and the transition state
found unrestricted can lie far from the one found restricted. A path
relaxed restricted then no longer covers the region the transition
state lies in when calculated unrestricted, and RQ2 becomes a real
question. OMol25 is exactly this case: every label recomputed
unrestricted at \mbox{$\omega$B97M-V/def2-TZVPD}, on the inherited
Transition1x geometries and on many structures beyond
them~\cite{omol25}. Chapter~\ref{ch:mr} asks RQ2 there.

RQ1 was the task this thesis was set. RQ2 arose during that work,
when the OMol25 models, used along the way as a point of comparison,
delivered structures other than the restricted reference transition
states this work had built and used for Chapter~\ref{ch:delta}, on
some reactions. Whether these models can find broken-symmetry
transition states at all had been tested nowhere, because the
published benchmarks only use restricted references, as the rest of
this chapter shows; the Preface gives the account.

The published benchmarks that test the OMol25 models at locating
transition states report high success rates. According to them, four
OMol25 models, three
of them used in Chapter~\ref{ch:mr}, reach the reference saddle in 91.8\,\% of
automated transition-state searches, and an MLIP search is recommended
as the default first stage of such searches~\cite{marks2026}. Hait et
al. use eSEN to produce the starting guess for a DFT transition-state
optimisation, and the DFT optimisation converges from that guess in
fewer steps than from a guess produced with DFT alone~\cite{hait2025}.

What these benchmarks do not measure is the outcome at broken-symmetry
transition states, because both test against transition states
optimised on the restricted surface. Hait et al. say so: they
re-optimised their 121 references unrestricted, 27 broke spin symmetry
and moved away from the supplied saddle, and all 27 were excluded as
defective references rather than replaced~\cite{hait2025}. The authors
expected models like eSEN to be challenged by these cases even with
unrestricted training data, and did not test it. Marks et al. report
nothing about the SCF solution of their references, so we ran the
stability analysis ourselves, at the paper's own level of theory. Of
the 25 even-electron references, 24 are closed-shell and one is
broken-symmetry (Appendix~\ref{app:marks}). That one was still optimised restricted, so it is not a transition
state of the unrestricted surface and should not serve as a
reference. It does, and so a model that finds the same restricted
structure is scored as a success, although the true transition state
lies elsewhere.

The structures are not rare. OMol25 reports $\langle S^2 \rangle >
0.001$, that is, broken-symmetry solutions, for just under 5\,\% of
its unrestricted Transition1x configurations~\cite{omol25}. That number counts every configuration of
a reaction path, most of which lie near reactant and product where the
restricted solution is stable, so the share at the saddle points
themselves is not known and can be assumed to be higher. In the
reference set of Hait et al. it is a fifth~\cite{hait2025}.

Where a transition state is an open-shell singlet, quantum chemistry
has a rule. The affordable treatment is broken-symmetry unrestricted
DFT~\cite{grafenstein2002}, and the geometry is optimised on that
surface, because the restricted saddle is not a stationary point of
it~\cite{crawford2001}. The label should be computed on the surface the
geometry was optimised on. The
benchmarks above and the training data behind the OMol25 models each
depart from this rule, in different ways (Table~\ref{tab:surfaces}).

\begin{table}[h]
\centering
\small
\caption{On which surface the geometry and the label of a
broken-symmetry transition state lie.}
\label{tab:surfaces}
\begin{tabular}{@{}lll@{}}
\toprule
 & geometry & label \\
\midrule
rule & broken-symmetry & broken-symmetry \\
benchmark references (Appendix~\ref{app:marks}, \cite{hait2025, marks2026}) & restricted & restricted \\
OMol25 training data~\cite{omol25} & restricted & broken-symmetry \\
\bottomrule
\end{tabular}
\end{table}

The benchmark references are restricted throughout, geometry and label,
which is why a model that reproduces the restricted saddle matches them
(Table~\ref{tab:surfaces}).
The OMol25 training data keeps half of the rule: the Transition1x
geometries come from restricted searches that never tested the
solution~\cite{schreiner2022}, and OMol25 recomputed their labels
unrestricted and left the geometries where they were~\cite{omol25}.
Whether half is enough has not been measured.

Chapter~\ref{ch:mr} attempts to close that gap. From the Transition1x
test split it selects reactions of high multireference character,
where broken-symmetry solutions are expected at the transition state,
together with a mid- and a low-MR group as control. The OMol25 models
UMA-S, UMA-M and eSEN then drive the NEB searches, and every transition
state they find is re-evaluated with unrestricted DFT, the protocol
OMol25 used for its labels, extended by a stability analysis. The
stability analysis decides whether the lowest solution at the structure
the model found is closed-shell or broken-symmetry, and it splits the
transition states into these two groups. On the two groups the same
metrics are applied, to test whether the models are equally successful
at finding a transition state in a broken-symmetry region as in a
closed-shell one.

They are not. At broken-symmetry transition states the models deliver
structures that are not stationary points of the unrestricted surface,
and their forces there deviate more from DFT than at closed-shell ones.
The chapter traces this to the training data. The training geometries of the broken-symmetry reactions are not
stationary points of the unrestricted surface: the force there is
not zero, so the models have never seen a transition state of that
surface, only the region around it.
